\section{Radical support and the exact conjecture}\label{sec:foundations}
\subsection{The multiplicative quantity that ignores repetition}
Every positive integer has a unique prime factorization
\[
 n=\prod_{p\mid n}p^{e_p},\qquad e_p\ge1.
\]
Its radical is
\[
 \rad(n)=\prod_{p\mid n}p,\qquad \rad(1)=1.
\]
The radical retains which primes occur and discards how many times they occur. Thus $\rad(72)=\rad(2^3\cdot3^2)=6$, whereas $72/6=12$ records the repeated-power contribution. If $n$ is squarefree, then $\rad(n)=n$.

For a positive primitive additive triple
\[
 a+b=c,\qquad \gcd(a,b)=1,
\]
all three pairs are coprime. For example, a divisor of $a$ and $c$ also divides $c-a=b$, hence is one. Therefore the prime supports are disjoint and
\[
 R=\rad(abc)=R_aR_bR_c,\qquad R_u=\rad(u).
\]
Put $M_u=u/R_u$. The exact decomposition is
\[
 abc=R M_aM_bM_c.
\]
The conjecture compares $c$ with $R$, rather than with $abc$ itself. Since $abc$ already contains $c$, an estimate against the full product would lose the intended content.

\figplaceholder{01_radical}{Prime support and repeated depth are different quantities. The radical keeps one copy of each active prime. The powerful part keeps the remaining copies.}

\begin{conjecture}[Masser--Oesterl\'e $abc$]\label{conj:abc}
For every real $\eps>0$ there is a constant $K_\eps>0$, depending only on $\eps$, such that every positive integer triple with $a+b=c$ and $\gcd(a,b)=1$ satisfies
\[
 \boxed{c\le K_\eps\,\rad(abc)^{1+\eps}.}
\]
\end{conjecture}
The order of the quantifiers is the statement. First choose $\eps$. Then choose a constant that works for every primitive triple. The constant may become large as $\eps$ approaches zero. It may not depend on the particular triple.

This formulation is the one stated in Stewart and Yu's primary treatment of the conjecture~\cite{StewartYu2001}. Replacing $\le$ by $<$ does not change the all-$\eps$ conjecture after enlarging the constant. Requiring $a\le b$ also leaves it unchanged because the two inputs can be exchanged.

\subsection{Equivalent output-sector formulations}
The form best adapted to this research separates the repeated depth of the output from the input prime support.
\begin{proposition}[Exact normalization]\label{prop:defect}
For every primitive positive triple,
\[
 \frac cR=\frac{M_c}{R_aR_b},\qquad
 Q:=\frac{abc}{R},\qquad \frac Q{ab}=\frac cR.
\]
Consequently Conjecture~\ref{conj:abc} is exactly equivalent to each all-$\eps$ assertion
\[
 M_c\le K_\eps R_aR_bR^\eps,
 \qquad Q\le K_\eps abR^\eps.
\]
\end{proposition}
\begin{proof}
Use $R=R_aR_bR_c$ and $c=R_cM_c$, then cancel the indicated factors. Multiplying any displayed inequality by the positive cancelled factors recovers the original conjecture, with the same constant and exponent.
\end{proof}
A local construction may measure $M_c$ naturally. A Wronskian may instead be divisible by $Q$. Both are legitimate, provided the final global height includes the matching normalization.

\subsection{Quality measures the exponent in one example}
Define the quality of a primitive triple by
\[
 q(a,b,c)=\frac{\log c}{\log R}.
\]
The denominator is positive because $c\ge2$ and $R\ge2$. The logarithm base does not matter. The identity $c=R^q$ is exact, although decimal qualities are approximations.

A triple has $q>1$ exactly when $c>R$. It then has more output size than the global distinct-prime product. Such a triple is often called an $abc$ hit. A hit is an interesting example, not a counterexample to the conjecture. The constant $K_\eps$ is part of the conjecture.

\begin{center}
\begin{tabular}{@{}rrrrl@{}}
\toprule
$a$&$b$&$c$&$R$&Factorization of the repeated terms\\
\midrule
$1$&$8$&$9$&$6$&$8=2^3$, $9=3^2$\\
$3$&$125$&$128$&$30$&$125=5^3$, $128=2^7$\\
$1$&$4374$&$4375$&$210$&$4374=2\cdot3^7$, $4375=5^4\cdot7$\\
$2$&$109\cdot3^{10}$&$23^5$&$15042$&$23^5=6436343$\\
\bottomrule
\end{tabular}
\end{center}
The last equality is $2+109\cdot3^{10}=23^5$. Every factor in the table matters. Its radical is $2\cdot3\cdot109\cdot23=15042$, giving quality about $1.629912$. The third triple has quality about $1.567887$. These are illustrative exact examples, without any claim that either is the largest known quality.

\figplaceholder{03_high_quality}{Exact examples compare the output with the radical. Decimal quality is a diagnostic derived from the exact integer data. A finite collection of high qualities does not decide the all-triple conjecture.}

\begin{proposition}[Finiteness formulation]
Conjecture~\ref{conj:abc} is equivalent to the assertion that, for every $\eps>0$, there are only finitely many primitive positive triples with $q(a,b,c)>1+\eps$.
\end{proposition}
\begin{proof}
Assume $abc$, using its bound with $\eps/2$. A triple of quality above $1+\eps$ satisfies
\[
 R^{1+\eps}<c\le K_{\eps/2}R^{1+\eps/2}.
\]
Hence $R^{\eps/2}<K_{\eps/2}$. The radical is bounded. The same conjectural inequality then bounds $c$, leaving finitely many positive triples.

Conversely, fix $\eps>0$ and assume the finiteness assertion. Outside its finite exceptional set, $c\le R^{1+\eps}$. Choose $K_\eps$ to be at least one and at least every ratio $c/R^{1+\eps}$ in the exceptional set. This single constant works for all triples.
\end{proof}
This equivalence explains why large searches can be valuable and still be logically limited. A computation can list hits in a finite range. It cannot certify the finiteness of the unsearched set.

\subsection{Why the epsilon cannot be removed}
The uniform statement $c\le KR$ for one absolute constant $K$ is false. There is an elementary family that proves this without using $abc$.
\begin{theorem}[Failure of the zero-epsilon estimate]\label{thm:epsilon-zero}
The ratios $c/\rad(abc)$ are unbounded among primitive positive additive triples.
\end{theorem}
\begin{proof}
For $n\ge0$, take
\[
 a_n=1,\quad b_n=4^{3^n}-1,\quad c_n=4^{3^n}.
\]
These are primitive, $b_n$ is odd, and $\rad(c_n)=2$. We claim
\[
 v_3(b_n)=n+1.
\]
At $n=0$ this is $v_3(3)=1$. If $x=1+3^ku$ with $k\ge1$ and $3\nmid u$, then
\[
 x^3-1=3^{k+1}u\left(1+3^ku+3^{2k-1}u^2\right).
\]
The bracket is a unit modulo three, so cubing increases the valuation by exactly one. This proves the claim by induction.

Because $b_n$ contains $3^{n+1}$, its powerful part contains $3^n$. Thus
\[
 \rad(b_n)\le \frac{b_n}{3^n},\qquad
 R_n=2\rad(b_n)\le\frac{2b_n}{3^n}.
\]
Therefore
\[
 \frac{c_n}{R_n}\ge\frac{3^n c_n}{2b_n}>\frac{3^n}{2}\longrightarrow\infty.
\]
\end{proof}
The displayed lower bound does not contradict the positive-$\eps$ conjecture. It does not exhibit infinitely many triples with a fixed positive quality gap. The output grows doubly exponentially with $n$, while the proved defect lower bound grows exponentially. A large defect lower bound and a fixed positive exponent gap are different assertions.

\subsection{A known unconditional bound}
Stewart and Yu proved an effective bound of a different scale~\cite[Theorem 1]{StewartYu2001}: there is an effectively computable absolute $C>0$ such that every primitive positive triple satisfies
\[
 c<\exp\!\left(CR^{1/3}(\log R)^3\right).
\]
It controls the height in terms of the radical unconditionally. Its exponential dependence is much larger than the power $R^{1+\eps}$ required by $abc$. It therefore provides a benchmark, without closing the gap this research addresses.

\subsection{A published claim and a continuing verification dispute}
Mochizuki's four IUT papers were published in 2021. The fourth paper states an application to $abc$~\cite{MochizukiIUTIV}. Publication and establishment of the claimed proof are distinct questions. Scholze and Stix documented an objection to the comparison underlying IUT III, Corollary 3.12, and Mochizuki published a response disputing their simplifications~\cite{ScholzeStix2018,MochizukiComments2018}.

A 2026 Project LANA verification report withholds a verdict on whether IUT proves $abc$. It identifies a needed compatibility between the relevant analytic and anabelian comparison maps~\cite{LANA2026}. That report also distinguishes its objection from accepting every identification in the earlier criticism. This edition follows that precise status: a published proof claim and a documented, unresolved verification dispute. None of the project proofs below depends on treating IUT as an established proof of $abc$, and none claims to repair its disputed comparison.
